\section{Related work}\label{sec:related}

\paragraph{Sparse moment--SOS relaxations.}
The moment--SOS hierarchy~\cite{lasserre2001global,parrilo2003semidefinite}
becomes impractical for long trajectories unless sparsity is exploited.
Correlative sparsity~\cite{waki2006sums,lasserre2006convergent} restricts
monomials to cliques of interacting variables, and term
sparsity~\cite{wang2021tssos,wang2022cstssos} further prunes the monomial bases.
Occupation-measure relaxations of optimal control~\cite{lasserre2008nonlinear}
and recent contact-rich planning relaxations~\cite{graesdal2024tight,kang2025spot}
give the trajectory setting that motivates this work. These papers reduce the
size of the relaxation. We take a fixed sparse relaxation as given and study the
cost of solving its IPM Newton systems.

\paragraph{Certifiable trajectory optimization solvers.}
STROM~\cite{kang2024strom} builds sparse second-order Lasserre relaxations of
polynomial trajectory problems and solves them with cuADMM, a symmetric
Gauss--Seidel ADMM implemented in CUDA, reporting certified suboptimality below
1\% on its benchmarks. TinySDP~\cite{mahajan2026tinysdp} combines PSD-cone
projections with a cached-Riccati ADMM for real-time model predictive control
with nonconvex obstacle constraints on microcontrollers. Both are first-order
methods: each iteration is cheap and parallel, but high accuracy takes many
iterations. MIST is complementary. It keeps the second-order interior-point
method, which reaches high accuracy in a few tens of iterations
(\cref{sec:pushbox}), and makes each Newton step cheap through exact temporal
elimination. On the local side, interior-point solvers for the nonlinear
trajectory problem itself, such as CALIPSO~\cite{howell2022calipso} for conic
and complementarity constraints and Hippo~\cite{zhao2026hippo}, can serve as
the refinement step that turns an extracted relaxation solution into a feasible
trajectory.

\paragraph{Chordal structure and the complexity of sparse SDPs.}
Chordal conversion and matrix completion~\cite{fukuda2001exploiting,nakata2003exploiting,kim2011exploiting,vandenberghe2015chordal}
split large PSD constraints into overlapping small ones, and
first-order methods exploit the same decomposition~\cite{zheng2020chordal}.
\citet{zhang2020clique} showed that dualized clique-tree conversion gives
near-linear-time IPMs for SDPs whose aggregate sparsity graph has small
treewidth. \citet{zhang2024chordal} showed that sparsity of the aggregate graph
alone does not control the Newton-system cost, because the incidence of
constraints on blocks must also be accounted for. Our analysis follows this line
but works directly on the coefficient-matching incidence of a fixed SOS
relaxation. It derives separators from shared monomials and certifies the
hierarchy used by the solver on the assembled Schur graph.

\paragraph{Structured interior-point methods for optimal control.}
For quadratic and nonlinear programs in optimal control, the stagewise structure
of the KKT system is classical: Riccati recursions and block-tridiagonal
factorizations give per-iteration cost linear in the
horizon~\cite{wright1991partitioned,rao1998application,frison2020hpipm}.
In SOS relaxations, the Newton variables are multipliers of coefficient
equations indexed by monomials, not states and controls. Temporal coupling
arises through monomials shared by PSD blocks, and pruning or coarsening can
break a naive stage structure. This is why we need separator bounds and a
certificate before a chain solver can be applied.

\paragraph{Sparse direct methods.}
The relation between elimination orderings and fill is
classical~\cite{rose1976algorithmic,liu1990role}, as are nested
dissection~\cite{george1973nested}, minimum-degree
orderings~\cite{amestoy1996approximate}, multifrontal and supernodal
factorizations~\cite{duff1983multifrontal,chen2008cholmod}, and their
surveys~\cite{davis2016survey}. Block cyclic reduction solves block-tridiagonal
systems with logarithmic depth~\cite{hockney1965fast,heller1976some}. We show
in \cref{prop:ordering} that our hierarchy is one particular elimination
ordering. Its value therefore lies in obtaining a certified ordering and an
efficient GPU execution plan from the model, not in a new factorization
algorithm. NVIDIA cuDSS~\cite{cudss} is the general sparse GPU baseline in our
experiments.

\paragraph{SDP solvers.}
Primal--dual IPMs for SDP~\cite{helmberg1996interior,nesterov1997self,todd1998nesterov},
with Mehrotra-type predictor--corrector steps~\cite{mehrotra1992implementation}
and homogeneous self-dual embeddings~\cite{ye1994homogeneous}, underlie
SeDuMi~\cite{sturm1999sedumi}, SDPT3~\cite{toh1999sdpt3} and
MOSEK~\cite{mosek}. \citet{papp2019sos} optimize directly over the SOS cone
with a nonsymmetric IPM in an interpolant basis, avoiding the SDP reformulation,
and \citet{leijenhorst2024clustered} develop a high-precision primal--dual IPM
for clustered low-rank SDPs arising from polynomial optimization. Parallel and GPU implementations have targeted the dense
Schur complement of general SDPs~\cite{fujisawa2012sdpara}. Low-rank and
first-order methods~\cite{burer2003nonlinear} scale to larger problems but
typically reach lower accuracy, and BRISK~\cite{brisk} is the second CPU
baseline in our experiments. \todo{one sentence describing BRISK's algorithm
and its reference.} On the theoretical side, fast IPMs for SDP and SOS maintain
approximate inverses of the Hessian across
iterations~\cite{jiang2020faster,jiang2022sos}. That work exploits reuse across
iterations, whereas we exploit temporal structure within one iteration;
\cref{app:inverse} discusses how the two might be combined.
